% !TeX program = pdflatex
\documentclass[11pt,a4paper]{amsart}

\newcommand{\notelanguage}{finnish}
% Shared preamble for course notes and exercise sets.

\usepackage[T1]{fontenc}
\usepackage{lmodern}
\providecommand{\notelanguage}{english}
\usepackage[finnish,english]{babel}
\AtBeginDocument{\expandafter\selectlanguage\expandafter{\notelanguage}}

\newcommand{\notesfinnishlanguage}{finnish}
\ifx\notelanguage\notesfinnishlanguage
\newcommand{\notestheoremname}{Lause}
\newcommand{\noteslemmaname}{Lemma}
\newcommand{\notespropositionname}{Propositio}
\newcommand{\notescorollaryname}{Seurauslause}
\newcommand{\notesdefinitionname}{Määritelmä}
\newcommand{\notesexamplename}{Esimerkki}
\newcommand{\notesproblemname}{Tehtävä}
\newcommand{\notesexercisename}{Tehtävä}
\newcommand{\notesremarkname}{Huomautus}
\newcommand{\notessolutionname}{Ratkaisu}
\else
\newcommand{\notestheoremname}{Theorem}
\newcommand{\noteslemmaname}{Lemma}
\newcommand{\notespropositionname}{Proposition}
\newcommand{\notescorollaryname}{Corollary}
\newcommand{\notesdefinitionname}{Definition}
\newcommand{\notesexamplename}{Example}
\newcommand{\notesproblemname}{Problem}
\newcommand{\notesexercisename}{Exercise}
\newcommand{\notesremarkname}{Remark}
\newcommand{\notessolutionname}{Solution}
\fi

\usepackage[a4paper,margin=30mm]{geometry}
\usepackage{microtype}
\usepackage{graphicx}
\usepackage{enumitem}

\newcommand{\exerciseheading}[2]{%
  \par\noindent
  {\large\bfseries #1\par}
  \vspace{0.1em}
  \noindent{\large #2\par}
  \vspace{1.5em}
}

\usepackage{amsmath}
\usepackage{amssymb}
\usepackage{amsthm}
\usepackage{mathtools}
\usepackage{aliascnt}

% Linear systems: the argument gives the number of alignment pairs.
\newenvironment{linearsystem}[1]
  {\left\{\begin{alignedat}{#1}}
  {\end{alignedat}\right.}

\usepackage[hidelinks,pdfusetitle]{hyperref}

\numberwithin{equation}{section}
\setlist{itemsep=0.25em,topsep=0.5em}

\theoremstyle{plain}
\newtheorem{theorem}{\notestheoremname}[section]
\newaliascnt{lemma}{theorem}
\newtheorem{lemma}[lemma]{\noteslemmaname}
\aliascntresetthe{lemma}
\newaliascnt{proposition}{theorem}
\newtheorem{proposition}[proposition]{\notespropositionname}
\aliascntresetthe{proposition}
\newaliascnt{corollary}{theorem}
\newtheorem{corollary}[corollary]{\notescorollaryname}
\aliascntresetthe{corollary}

\theoremstyle{definition}
\newaliascnt{definition}{theorem}
\newtheorem{definition}[definition]{\notesdefinitionname}
\aliascntresetthe{definition}
\newaliascnt{example}{theorem}
\newtheorem{example}[example]{\notesexamplename}
\aliascntresetthe{example}
\newaliascnt{problem}{theorem}
\newtheorem{problem}[problem]{\notesproblemname}
\aliascntresetthe{problem}
\newtheorem{exercise}{\notesexercisename}

\theoremstyle{remark}
\newaliascnt{remark}{theorem}
\newtheorem{remark}[remark]{\notesremarkname}
\aliascntresetthe{remark}

\newenvironment{solution}
{
\begin{proof}[\notessolutionname]}
  {
\end{proof}}

\usepackage[nameinlink,noabbrev]{cleveref}

\crefname{theorem}{theorem}{theorems}
\crefname{lemma}{lemma}{lemmas}
\crefname{proposition}{proposition}{propositions}
\crefname{corollary}{corollary}{corollaries}
\crefname{definition}{definition}{definitions}
\crefname{example}{example}{examples}
\crefname{problem}{problem}{problems}
\crefname{exercise}{exercise}{exercises}
\crefname{remark}{remark}{remarks}

\ifx\notelanguage\notesfinnishlanguage
\crefname{theorem}{lause}{lauseet}
\crefname{lemma}{lemma}{lemmat}
\crefname{proposition}{propositio}{propositiot}
\crefname{corollary}{seurauslause}{seurauslauseet}
\crefname{definition}{määritelmä}{määritelmät}
\crefname{example}{esimerkki}{esimerkit}
\crefname{problem}{tehtävä}{tehtävät}
\crefname{exercise}{tehtävä}{tehtävät}
\crefname{remark}{huomautus}{huomautukset}
\fi

\newcommand{\NN}{\mathbb{N}}
\newcommand{\NNone}{\mathbb{N}_1}
\newcommand{\ZZ}{\mathbb{Z}}
\newcommand{\QQ}{\mathbb{Q}}
\newcommand{\RR}{\mathbb{R}}
\newcommand{\CC}{\mathbb{C}}
\newcommand{\PP}{\mathbb{P}}

\DeclareMathOperator{\im}{im}
\DeclareMathOperator{\rank}{rank}
\DeclareMathOperator{\syt}{syt}
\DeclarePairedDelimiter{\abs}{\lvert}{\rvert}
\DeclarePairedDelimiter{\norm}{\lVert}{\rVert}
\newcommand{\set}[1]{\left\lbrace#1\right\rbrace}
\newcommand{\maxset}[1]{\max\set{#1}}
\newcommand{\minset}[1]{\min\set{#1}}
\newcommand{\powerset}[1]{\mathcal{P}\!\left(#1\right)}

\counterwithout{equation}{section}
\counterwithout{theorem}{section}

\title{Analyysi II\\Tutkijapolku 1}
\author{}
\date{}

\begin{document}

\exerciseheading{Analyysi II}{Tutkijapolku 1}

\begin{lemma}\label{lem:cont-seq}
  Olkoon \(f\colon (a,b) \to \RR\) jatkuva pisteessä \(x \in (a,b)\).
  Jos \(x_n \in (a,b)\) kaikilla \(n \in \NNone\) ja \(x_n \to x\),
  niin \(f(x_n) \to f(x)\).
\end{lemma}

\begin{proof}
  Olkoon \(\varepsilon > 0\). Jatkuvuuden nojalla on olemassa
  \(\delta > 0\) siten, että \(\abs{f(t)-f(x)} < \varepsilon\), kun
  \(t \in (a,b)\) ja \(\abs{t-x} < \delta\).
  Koska \(x_n \to x\), on olemassa \(N \in \NNone\) siten, että
  \(\abs{x_n-x} < \delta\) kaikilla \(n \geq N\). Näin ollen
  \(\abs{f(x_n)-f(x)} < \varepsilon\) kaikilla \(n \geq N\),
  joten \(f(x_n) \to f(x)\).
\end{proof}

\begin{exercise}
  Määritä kaikki jatkuvat funktiot \(f\colon \RR \to \RR\), joille
  \begin{equation}\label{eq:add}
    f(x+y) = f(x) + f(y)
  \end{equation}
  kaikilla \(x,y \in \RR\). Mitä tapahtuu, jos ei oleteta jatkuvuutta?
\end{exercise}

\begin{solution}
  Ensiksi ehdosta \eqref{eq:add} seuraa, että
  \[
    f(0) = f(0 + 0) = 2f(0),
  \]
  eli \(f(0) = 0\).

  Olkoon \(x \in \RR\) mielivaltainen. Näytetään induktiolla, että
  \(f(nx) = nf(x)\) kaikilla \(n \in \NN\). Ensiksi
  \(f(0x) = 0 = 0f(x)\), joten väite pätee, kun \(n = 0\). Olkoon sitten
  \(n \in \NN\) mielivaltainen ja oletetaan, että
  \[
    f(nx) = nf(x).
  \]
  Tällöin yhtälön \eqref{eq:add} nojalla
  \[
    f((n + 1)x) = f(nx + x) = f(nx) + f(x) = nf(x) + f(x) = (n + 1)f(x).
  \]
  Väite pätee siis luvulla \(n + 1\). Täten väite pätee kaikilla \(n \in \NN\).

  Lisäksi
  \[
    0 = f(x + (-x)) = f(x) + f(-x),
  \]
  eli \(f(-x) = -f(x)\). Näistä yhdessä seuraa, että \(f(zx) = zf(x)\)
  kaikilla \(z \in \ZZ\) ja \(x \in \RR\). Erityisesti asettamalla
  \(c = f(1)\) saadaan \(f(z) = zf(1) = cz\) kaikilla \(z \in \ZZ\).

  Olkoon sitten \(m \in \ZZ\) ja \(n \in \NNone\). Tällöin
  \[
    nf\left(\frac{m}{n}\right) = f\left(n\frac{m}{n}\right) = f(m) = cm.
  \]
  Nyt jakamalla puolittain luvulla \(n\) saadaan
  \[
    f(q) = cq
  \]
  kaikilla \(q \in \QQ\).

  Olkoon \(x \in \RR\) mielivaltainen. Rationaalilukujen tiheyden nojalla
  on olemassa lukujono \((q_k)\), missä \(q_k \in \QQ\) kaikilla
  \(k \in \NNone\) ja \(q_k \to x\), kun \(k \to \infty\).
  Lemman \ref{lem:cont-seq} nojalla
  \[
    f(x) = \lim_{k \to \infty} f(q_k) = \lim_{k \to \infty} cq_k = cx.
  \]
  Olemme siis tässä vaiheessa osoittaneet, että kaikki jatkuvat ja
  ehdon \ref{eq:add} täyttävät funktiot ovat muotoa \(f(x) = cx\),
  missä \(c = f(1) \in \RR\).

  On siis vielä näytettävä, että funktiot \(f \colon \RR \to \RR\),
  \(f(x) = cx\),
  missä \(c \in \RR\) ovat jatkuvia ja toteuttavat ehdon \ref{eq:add}.

  Olkoon \(\varepsilon > 0\) ja \(t \in \RR\). Valitaan \(\delta =
  \frac{\varepsilon}{\abs{c} + 1} > 0\). Tällöin
  \[
    \begin{aligned}
      \abs{f(x) - f(t)} &= \abs{c(x - t)}  \\
      &= \abs{c}\abs{x - t} \\
      &\leq \frac{\varepsilon\abs{c}}{\abs{c} + 1} \\
      &< \varepsilon
    \end{aligned}
  \]
  kaikilla \(x \in \RR\) joilla \(\abs{x - t} < \delta\). Täten \(f\) on jatkuva
  pisteessä \(t \in \RR\). Koska \(t\) oli mielivaltainen, \(f\) on jatkuva.
  Lisäksi,
  \[
    f(x + y) = c(x + y) = cx + cy = f(x) + f(y)
  \]
  kaikilla \(x,y \in \RR\), joten \(f\) täyttää myös ehdon \ref{eq:add}.

  Olemme nyt osoittaneet, että jatkuvat, ehdon \ref{eq:add} täyttävät
  funktiot ovat täsmälleen reaalisen vektoriavaruuden \(\RR\)
  lineaariset operaattorit
  eli funktiot \(f \colon \RR \to \RR\), \(f(x) = cx\), missä \(c \in \RR\).

  Tutkitaan seuraavaksi ehdon \ref{eq:add} täyttävien funktioiden \(f\)
  käyttäytymistä ilman oletuksia jatkuvuudesta. Edellisen perusteella
  tiedetään, että ehdon \ref{eq:add} täyttävät funktiot ovat muotoa
  \[
    f(q) = cq
  \]
  kaikilla \(q \in \QQ\). Tämän laajentaminen reaaliluvuille käytti funktion
  \(f\) jatkuvuutta.

  Olkoon \(x \in \RR\), \(m \in \ZZ\) ja \(n \in \NNone\). Edellisen ja
  ehdon \ref{eq:add} nojalla
  \[
    nf\left(\frac{m}{n}x\right) = f(mx) = mf(x),
  \]
  joten \(f(qx) = qf(x)\) kaikilla \(q \in \QQ\) ja \(x \in \RR\).
  Täten additiivisuuden ja yhtälön \(f(qx) = qf(x)\) nojalla
  \(f\) on \(\QQ\)-lineaarinen.

  Olemme siis osoittaneet, että jokainen ehdon \eqref{eq:add} täyttävä
  funktio on \(\QQ\)-lineaarinen. Lisäksi jokainen \(\QQ\)-lineaarinen
  kuvaus on additiivinen, joten ilman jatkuvuusoletusta ratkaisut ovat
  täsmälleen \(\QQ\)-lineaariset kuvaukset \(\RR \to \RR\).

  Näytetään vielä, että tällainen kuvaus ei välttämättä ole jatkuva.
  Tarkastellaan reaalilukuja vektoriavaruutena rationaalilukujen yli.
  Koska \(\sqrt{2}\) on irrationaalinen, luvut \(1\) ja \(\sqrt{2}\)
  ovat lineaarisesti riippumattomia yli \(\QQ\):n. Kannanlaajennuslauseen
  nojalla joukko \(\set{1,\sqrt{2}}\) voidaan täydentää tämän
  vektoriavaruuden kannaksi \(\mathcal{B}\). Asetetaan
  \[
    f(\sqrt{2}) = 1, \qquad
    f(b) = 0 \quad \text{kaikilla } b \in \mathcal{B} \setminus \set{\sqrt{2}}.
  \]
  Kannalla määritelty kuvaus voidaan laajentaa yksikäsitteisesti
  \(\QQ\)-lineaariseksi kuvaukseksi \(f\colon \RR \to \RR\).
  Täten \(f\) täyttää ehdon \eqref{eq:add}.

  Jos \(f\) olisi jatkuva, aiemmin osoitetun perusteella olisi
  \[
    f(x) = f(1)x = 0
  \]
  kaikilla \(x \in \RR\). Tämä on ristiriidassa ehdon
  \(f(\sqrt{2}) = 1\) kanssa. Täten \(f\) ei ole jatkuva.
\end{solution}

\begin{remark}
  Edellä käytettiin ilman todistusta kannanlaajennuslausetta sekä sitä,
  että kannalla määritelty kuvaus voidaan laajentaa yksikäsitteisesti
  lineaariseksi kuvaukseksi.
\end{remark}

\end{document}
